\section{FreeRTOS Adoption: Design Rationale}

\subsection{Motivation}

The DAQ firmware must concurrently manage several distinct real-time responsibilities:
high-rate ADC sampling, PTP clock synchronisation (IEEE~1588), MQTT telemetry
publication, and peripheral management (Ethernet, SPI/I\textsuperscript{2}C).
Implementing these as a monolithic superloop introduces structural problems: timing
jitter between subsystems, polling overhead competing with acquisition, and brittle
sequencing logic that degrades maintainability as complexity grows.

FreeRTOS is selected as the firmware RTOS to address these concerns.

\subsection{Task Decomposition Over Superloop Architecture}

A bare-metal superloop executes all responsibilities sequentially within a single loop
iteration. Any blocking operation (waiting for an Ethernet ACK, a PTP timestamp
correction, or an MQTT broker handshake) directly stalls the entire loop, introducing
indeterminate latency across all subsystems.

FreeRTOS enables preemptive multitasking: each logical responsibility is encapsulated
as an independent task with its own stack and priority assignment.

\begin{table}[h]
    \centering
    \caption{FreeRTOS Task Allocation}
    \label{tab:freertos-tasks}
    \begin{tabular}{lll}
      \hline
        \textbf{Task} & \textbf{Priority} & \textbf{Responsibility} \\
        \hline
        \texttt{ADC\_SampleTask}    & High   & Triggered acquisition, DMA management \\
        \texttt{PTP\_SyncTask}      & High   & IEEE~1588 clock discipline \\
        \texttt{MQTT\_PublishTask}  & Normal & Packetising and publishing telemetry \\
        \hline
    \end{tabular}
\end{table}

This decomposition maps cleanly onto the system's logical architecture and makes the
firmware structure self-documenting. Inter-task dependencies are expressed explicitly
through FreeRTOS queue and semaphore APIs, rather than implicitly through shared
globals and flag polling.

\subsection{ADC Sampling Integrity}

A common concern when introducing an RTOS is scheduler-induced jitter disrupting
deterministic sampling. This does not apply here because ADC acquisition is decoupled
from the scheduler via DMA and hardware timers.

The STM32 ADC is configured in timer-triggered DMA mode:

\begin{itemize}[noitemsep]
    \item A hardware timer (e.g.\ TIM2) fires at a precise, crystal-referenced
          interval; entirely independent of the FreeRTOS tick.
    \item Each timer event triggers an ADC conversion; the result is transferred to a
          memory buffer by DMA with zero CPU involvement.
    \item \texttt{ADC\_SampleTask} wakes only when DMA signals a buffer-full condition
          (via \texttt{xTaskNotifyFromISR} or a binary semaphore), consuming completed
          samples rather than performing acquisition itself.
\end{itemize}

The scheduler therefore has no influence over the ADC trigger cadence. Sampling jitter
is determined by the hardware timer and ADC clock, not by task preemption. 

\subsection{Development Velocity and Ecosystem Maturity}

FreeRTOS is the dominant RTOS in the STM32 ecosystem. ST's HAL and CubeMX toolchain
provide first-party FreeRTOS integration via the CMSIS-RTOS~v2 wrapper, including
pre-validated configurations for the STM32 device family. Direct project benefits
include:

\begin{itemize}[noitemsep]
    \item \textbf{Synchronisation primitives}: mutexes, queues, and semaphores cover
          all inter-task communication needs (e.g.\ passing sample buffers from
          \texttt{ADC\_SampleTask} to \texttt{MQTT\_PublishTask}) without bespoke
          implementation.
    \item \textbf{Network stack integration}: LwIP, the embedded TCP/IP stack used
          for MQTT over Ethernet, has a documented FreeRTOS threading model with
          known-good integration patterns.
    \item \textbf{PTP driver compatibility}: IEEE~1588 hardware timestamp units on
          STM32 are accessed via HAL; the FreeRTOS task model integrates naturally with
          the interrupt-driven PTP servo loop.
    \item \textbf{Reduced driver development burden}: ST-provided middleware
          (ETH HAL, LwIP, FreeRTOS port) is tested and documented, avoiding hand-rolled
          scheduling, stack management, or network primitives.
\end{itemize}

This materially reduces implementation risk and accelerates development relative to a
bare-metal approach.

\subsection{Summary}

FreeRTOS is adopted because the DAQ node's responsibilities are genuinely concurrent,
and the RTOS provides the correct abstraction for managing that concurrency safely and
legibly. It introduces no risk to ADC sampling fidelity because acquisition is
hardware-driven and DMA-managed, making it scheduler-independent by design.
